\section{Discussion}

We developed a partial-view self-supervised framework for learning global representations of
DNA methylation samples. Unlike objectives that operate only at the level of masked features,
WCED routes reconstruction through a CLS bottleneck and evaluates reconstruction at CpGs not
provided to the encoder. Contrastive view matching directly trains the same sample embedding
that is subsequently transferred to downstream analysis. The combination targets a distinction
that is easy to miss in high-dimensional molecular modeling: successful prediction of masked
features is not equivalent to learning a useful global representation.

The current results establish three narrow but meaningful points. First, the encoder can
compress a partial methylation profile while retaining information about unobserved sites.
Second, independently sampled views of the same profile remain identifiable in representation
space. Third, the resulting encoder can be adapted reproducibly to age prediction. The CLS
representation contained age-associated information before fine-tuning and retained more
usable variation than mean pooling, consistent with the intended sample-level bottleneck.

The paired benchmark establishes better predictive performance for the implemented
\method{} pipeline than for MethylGPT on the matched held-out samples. It does not, however,
isolate the source of that difference: the pipelines vary in architecture, optimization, and
their handling of 1,760 CpGs outside the shared pretraining reference. ElasticNet provided a
strong non-neural reference, with performance close to the \method{} ensemble. Although all
three point estimates favored \method{}, sample-level ElasticNet predictions were not retained,
so statistical uncertainty is currently available only for the paired MethylGPT comparison.
Matched objective and initialization ablations are still required before attributing the
downstream gain specifically to WCED or to self-supervised pretraining.

The strongest limitation is cohort structure. Dataset identity is highly recoverable from the
pretrained embeddings, showing that study- or batch-specific signals remain prominent. Random
sample partitions can therefore overestimate transfer when samples from the same source study
appear in model development and evaluation. This issue also changes how representation
visualizations should be read. Separation by tissue may reflect genuine biology, but clean
clusters can coexist with strong study effects. Quantitative study-held-out tests should carry
more evidential weight than visually distinct UMAP clusters.

The representation analyses are therefore best viewed as diagnostics. Attention concentration,
genomic locality, and CpG embedding structure can motivate hypotheses, but they do not by
themselves identify causal age-associated loci or establish biological mechanism. Stronger
interpretation would require stability across folds and seeds, perturbation-based importance,
enrichment against a prespecified CpG background, and correction for multiple testing.

Further limitations include reliance on array-selected CpGs, incomplete evaluation across
platforms, a single firmly established downstream phenotype, and no prospective clinical
validation. A broad foundation-model claim would require transfer across multiple tasks and
cohorts. At the present stage, ``self-supervised methylation encoder'' or ``methylation
representation model'' more precisely describes the demonstrated scope.

In summary, the evidence supports WCED as a method for learning partial-view-consistent global
methylation representations and supports \method{} as a stable age-prediction encoder on the
held-out evaluation set. Whether the objective improves upon matched alternatives and whether
the representation transfers to unseen studies remain the decisive experiments for the final
manuscript.

% Internal drafting note (not printed): Human review: after the controlled comparisons are
%   complete, rewrite the opening and closing Discussion paragraphs around one
%   experimentally attributable conclusion. If no significant pretraining benefit is found,
%   present the paper as an objective evaluation of sample-level bottleneck learning rather
%   than forcing a superiority narrative.
